%========================
% Journal Article Template
%========================
\documentclass[10pt,a4paper]{article}

%---- Packages dasar
\usepackage[a4paper,margin=1in]{geometry}
\usepackage{graphicx}
\usepackage{booktabs}
\usepackage{siunitx}
\usepackage{amsmath,amssymb}
\usepackage{authblk}           % untuk afiliasi penulis
\usepackage[hidelinks]{hyperref}
\usepackage[capitalise,noabbrev]{cleveref}
\usepackage{caption}
\usepackage{subcaption}
\usepackage{setspace}
\usepackage[protrusion=true,expansion=true,tracking=true,kerning=true,spacing=true,activate={true,nocompatibility},final,disable=footnote]{microtype}
\usepackage[english, indonesian]{babel} % dukung EN & ID
\usepackage[backend=biber,style=ieee,sorting=none]{biblatex} % gaya IEEE
\addbibresource{references.bib} % ganti sesuai nama .bib kamu

%---- Info artikel
\title{\textbf{Judul Artikel Ilmiah yang Informatif dan Ringkas}}
\author[1]{Muhammad Reesa Rosyid}
\author[1,2]{Lubna Mawaddah}
\affil[1]{Faculty of Computer Science, Universitas Dian Nuswantoro, Semarang, Indonesia}
\affil[2]{Institusi Kedua (opsional)}
\date{\today}

\begin{document}
	\maketitle
	
	%========================
	% ABSTRACT (ENGLISH)
	%========================
	\selectlanguage{english}
	\begin{abstract}
		Write the English abstract here (150–250 words): background, objective, methods, key quantitative results, and implications.
	\end{abstract}
	\noindent\textbf{Keywords—} keyword 1; keyword 2; keyword 3.
	
	%========================
	% MAIN TEXT
	%========================
	\selectlanguage{english} % atau \selectlanguage{indonesian} untuk naskah penuh ID
	
	\section{Introduction}
	%Explain "what is pore" and its functionality of on real world application.
	
	Pori didefinisikan sebagai rongga atau ruang kosong di dalam matriks padat yang mampu menampung fluida serta memungkinkan pergerakannya melalui jaringan yang saling terhubung \cite{yi2022, wang2022}. Karakterisasi yang tepat serta pemahaman mendalam terhadap sifat-sifat geometris dari struktur pori memiliki pengaruh signifikan terhadap berbagai bidang ilmu, terutama dalam rekayasa dan geosains \cite{lee2022}. Sebagai contoh, dalam rekayasa perminyakan dan hidrologi, analisis struktur pori digunakan untuk menilai kualitas reservoir dan efisiensi perolehan fluida, sedangkan dalam rekayasa sipil dan ilmu material, pendekatan serupa dimanfaatkan untuk mengoptimalkan kekuatan serta daya tahan material berpori \cite{chen2024, xiao2023, liu2020, hou2023, liu2022}. Media berpori sendiri umumnya diklasifikasikan berdasarkan tingkat porositas, permeabilitas, serta sifat matriks padatnya, dengan variasi struktur yang dapat berkisar dari konfigurasi teratur pada material sintetis hingga bentuk yang sangat heterogen sebagaimana dijumpai pada tanah alami dan batuan retak \cite{yang2024}. Kompleksitas struktur tersebut menuntut penerapan pendekatan multi-skala dalam mempelajari fenomena aliran dan transportasi fluida, yang mencakup interaksi pada skala pori hingga proses yang berlangsung pada skala lapangan.
	
	
	%Why pores rock analysis are important in geophysics area, explain traditional methods and the limitations.
	
	%Explain about implementation with traditional ML and DL technique follow with prior studies and declare out of our research.
	
	%Explain general step of our method.
	
	\section{Methods}
	Uraikan metodologi secara replikatif: data, praproses, arsitektur/model, metrik evaluasi, dan protokol eksperimen.
	\subsection{Dataset and Preprocessing}
	Deskripsikan sumber data, ukuran, pembagian train/val/test, serta pembersihan/augmentasi.
	\subsection{Model / Algorithm}
	Detail komponen inti, rumus penting:
	\begin{equation}
		\mathcal{L} = \frac{1}{N}\sum_{i=1}^{N} \ell\big(f_\theta(x_i), y_i\big)
	\end{equation}
	\subsection{Evaluation Metrics}
	Definisikan metrik (Accuracy, Precision, Recall, F1, AUC, IoU, dsb.) dan alasan pemilihan.
	
	\section{Results}
	Paparkan hasil utama secara kuantitatif. Gunakan tabel/gambar yang jelas.
	\subsection{Quantitative Results}
	\begin{table}[h]
		\centering
		\caption{Perbandingan kinerja model pada set uji.}
		\label{tab:results}
		\begin{tabular}{lcccc}
			\toprule
			Model & Accuracy & Precision & Recall & F1 \\
			\midrule
			Baseline & 0.832 & 0.84 & 0.81 & 0.82 \\
			Proposed & \textbf{0.861} & \textbf{0.87} & \textbf{0.84} & \textbf{0.85} \\
			\bottomrule
		\end{tabular}
	\end{table}
	
	%\begin{figure}[h]
	%\centering
	%\includegraphics[width=.8\linewidth]{figures/sample_result.pdf}
	%\caption{Contoh hasil (beri keterangan singkat dan informatif).}
	%\label{fig:sample}
	%\end{figure}
	
	\section{Discussion}
	Interpretasikan hasil, jelaskan mengapa pendekatan bekerja/tidak bekerja, analisis error/ablation, dan keterbatasan studi.
	
	\section{Conclusion and Future Work}
	Ringkas kontribusi dan temuan utama. Nyatakan rencana pekerjaan lanjutan yang spesifik.
	
	\section*{Acknowledgements}
	Ucapkan terima kasih kepada pendanaan, kolaborator, atau fasilitas komputasi (tanpa klaim dukungan yang tidak benar).
	
	%========================
	% REFERENSI
	%========================
	% Contoh kutipan di teks: \cite{example2024paper}
	\printbibliography
	
	%========================
	% LAMPIRAN (opsional)
	%========================
	\appendix
	\section{Additional Experiments}
	Detail tambahan, konfigurasi hyperparameter, atau bukti matematis.
	
\end{document}
